\section{Compromiso ético y proyección comunitaria}
\label{sec:compromiso_comunitario}

Las autoras asumen la tecnología no como un fin lucrativo en sí mismo, sino como un instrumento indispensable de emancipación, justicia social y servicio comunitario. San Julián es una región con inmenso potencial económico que requiere la digitalización de sus instituciones educativas, agroproductivas y gremiales.

En este marco, se asume el compromiso ético de realizar talleres voluntarios y gratuitos de alfabetización digital y uso de herramientas web en favor de las nuevas promociones de secundaria del módulo. Asimismo, se contempla la posibilidad futura de conformar un emprendimiento juvenil de servicios informáticos y consultoría de software para acompañar los procesos de modernización tecnológica en el norte integrado cruceño.
